\section{Introduction}

A differential filter processes two input voltages before forming the output difference. Unequal processing converts part of a common input into differential output. \citet{casas1996basics} described matching the transfer functions of noncoupled paths and coupling single-ended filters through a floating common point as two approaches to common-mode rejection. The grounded ladders considered here follow the noncoupled-path configuration.

The relevant matching quantities depend on the measurement frequency and the connected load. \citet{hornero2021cm} showed that component mismatch can produce gain peaking and positive phase shifts in the common-mode response. For a passive ladder, the connected load also draws current through the output impedance. A component combination selected from the open-circuit voltage transfer therefore requires an output-impedance description for predicting the loaded response.

In a buffered resistive-sensor interface, the sample voltage is proportional to $I(t)R(t)$ and the reference voltage to $I(t)R_s$. When resistance changes during the observation window, filtering before forming the ratio produces a weighted history dependence. With mismatched channel responses, changing current can also affect a constant-resistance reading. A matching specification must therefore identify both the terminal conditions for the linear network and the dynamic quantity that the ratio readout estimates.

The objective of this study is to quantify the component freedom available for matching multistage grounded ladders under specified terminations and connect the matching criteria to dynamic ratio measurements. We used classical ladder extraction to express positive realizations in coefficient coordinates, then classified raw and calibrated matching under terminal constraints. Exact constructions and ideal-element circuit simulations checked the load dependence. A convolution analysis separated constant-resistance excitation rejection, resistance-history weighting and the first-order effect of channel-delay mismatch.

\begin{figure}[htbp]
\centering
\includegraphics[width=\linewidth]{figures/Fig_11_en.pdf}
\caption{Signal flow for the buffered ratio model. The highlighted block is the passive network analyzed by the coefficient equations; loads and probe admittances enter at the network output plane. The sensor changes the upstream source amplitude.}
\label{fig:system-plane}
\end{figure}

